\textbf{Phần 1. Trắc nghiệm lựa chọn}
\setcounterpageref{ex}{0}

\begin{ex} %[3D3V32]
Điền số thích hợp vào chỗ chấm: $1\ell = \dots\text{ ml}$.
\choice
{$100\text{ ml}$}
{\True $1000\text{ ml}$}
{$10\text{ ml}$}
{$500\text{ ml}$}
\loigiai{
$1\ell = 1000\text{ ml}$.
\textbf{Đáp án của câu hỏi là B}
}
\end{ex}

\begin{ex} %[3D3V32]
Một bình nước có dung tích $1\ell$. Bạn Nam đã uống hết $350\text{ ml}$, bạn Hùng uống hết $250\text{ ml}$. Trong bình còn lại bao nhiêu mi-li-lít nước?
\choice
{$500\text{ ml}$}
{$600\text{ ml}$}
{$300\text{ ml}$}
{\True $400\text{ ml}$}
\loigiai{
Đổi $1\ell = 1000\text{ ml}$. Cả hai bạn uống: $350 + 250 = 600\text{ ml}$. Còn lại: $1000 - 600 = 400\text{ (ml)}$.
\textbf{Đáp án của câu hỏi là D}
}
\end{ex}

\begin{ex} %[3D3V32]
Một lọ si-rô ho chứa $120\text{ ml}$. Mỗi ngày em bé uống 3 lần, mỗi lần $5\text{ ml}$. Hỏi lọ si-rô đó đủ cho em bé uống trong bao nhiêu ngày?
\choice
{\True $8$ ngày}
{$6$ ngày}
{$10$ ngày}
{$12$ ngày}
\loigiai{
Mỗi ngày em bé uống: $5 \times 3 = 15\text{ ml}$. Lọ si-rô uống trong: $120 : 15 = 8$ (ngày).
\textbf{Đáp án của câu hỏi là A}
}
\end{ex}

\begin{ex} %[3D3V32]
Có 6 chai nước khoáng, mỗi chai chứa $500\text{ ml}$. Tổng dung tích nước của cả 6 chai là
\choice
{$300\text{ ml}$}
{$3000\ell$}
{\True $3\ell$}
{$2\ell$}
\loigiai{
Tổng dung tích là: $500 \times 6 = 3000\text{ ml} = 3\ell$.
\textbf{Đáp án của câu hỏi là C}
}
\end{ex}

\begin{ex} %[3D3V32]
Dùng một cốc chia vạch loại $200\text{ ml}$ đong nước vào một chiếc ca dung tích $1\ell$. Hỏi cần đong ít nhất bao nhiêu cốc đầy để đầy chiếc ca đó?
\choice
{$4$ cốc}
{\True $5$ cốc}
{$6$ cốc}
{$10$ cốc}
\loigiai{
Đổi $1\ell = 1000\text{ ml}$. Số cốc cần là: $1000 : 200 = 5$ cốc.
\textbf{Đáp án của câu hỏi là B}
}
\end{ex}

\begin{ex} %[3D3V32]
Một hộp sữa tươi có dung tích $180\text{ ml}$. Một lốc sữa gồm 4 hộp như thế có tổng dung tích sữa là
\choice
{$680\text{ ml}$}
{$740\text{ ml}$}
{$800\text{ ml}$}
{\True $720\text{ ml}$}
\loigiai{
Tổng dung tích sữa của 1 lốc là: $180 \times 4 = 720\text{ (ml)}$.
\textbf{Đáp án của câu hỏi là D}
}
\end{ex}

\begin{ex} %[3D3V32]
Tính giá trị của biểu thức: $450\text{ ml} + 150\text{ ml} \times 3$.
\choice
{\True $900\text{ ml}$}
{$800\text{ ml}$}
{$1000\text{ ml}$}
{$950\text{ ml}$}
\loigiai{
$150 \times 3 = 450\text{ ml}$; $450 + 450 = 900\text{ ml}$.
\textbf{Đáp án của câu hỏi là A}
}
\end{ex}

\begin{ex} %[3D3V32]
Một bình thủy chứa $2\ell$ nước sôi. Mẹ rót đều vào 4 chiếc ca nhỏ. Mỗi ca chứa bao nhiêu mi-li-lít nước?
\choice
{$400\text{ ml}$}
{$600\text{ ml}$}
{\True $500\text{ ml}$}
{$250\text{ ml}$}
\loigiai{
Đổi $2\ell = 2000\text{ ml}$. Mỗi ca chứa: $2000 : 4 = 500\text{ (ml)}$.
\textbf{Đáp án của câu hỏi là C}
}
\end{ex}

\begin{ex} %[3D3V32]
Một chai dầu ăn có $750\text{ ml}$. Sau 2 tuần nấu ăn, trong chai còn lại $150\text{ ml}$. Trung bình mỗi tuần gia đình đã dùng bao nhiêu mi-li-lít dầu ăn?
\choice
{$600\text{ ml}$}
{\True $300\text{ ml}$}
{$250\text{ ml}$}
{$350\text{ ml}$}
\loigiai{
Số dầu đã dùng trong 2 tuần: $750 - 150 = 600\text{ ml}$. Mỗi tuần dùng: $600 : 2 = 300\text{ (ml)}$.
\textbf{Đáp án của câu hỏi là B}
}
\end{ex}

\begin{ex} %[3D3V32]
Để pha một ly nước cam ép, bạn An dùng $80\text{ ml}$ nước cốt cam và $120\text{ ml}$ nước lọc. Hỏi để pha 4 ly như thế cần bao nhiêu mi-li-lít nước ép cam thành phẩm?
\choice
{$600\text{ ml}$}
{$750\text{ ml}$}
{$900\text{ ml}$}
{\True $800\text{ ml}$}
\loigiai{
Một ly có dung tích: $80 + 120 = 200\text{ ml}$. Bốn ly cần: $200 \times 4 = 800\text{ (ml)}$.
\textbf{Đáp án của câu hỏi là D}
}
\end{ex}

\begin{ex} %[3D3V32]
So sánh: $A = 1\ell - 250\text{ ml}$ và $B = 350\text{ ml} \times 2$.
\choice
{\True $A > B$}
{$A < B$}
{$A = B$}
{Không so sánh được}
\loigiai{
$A = 1000 - 250 = 750\text{ ml}$. $B = 350 \times 2 = 700\text{ ml}$. Vì $750 > 700$ nên $A > B$.
\textbf{Đáp án của câu hỏi là A}
}
\end{ex}

\begin{ex} %[3D3V32]
Một bể cá cảnh nhỏ chứa $5\ell$ nước. Người ta dùng một chiếc ca dung tích $250\text{ ml}$ múc nước ra ngoài 8 lần. Trong bể còn lại bao nhiêu lít nước?
\choice
{$2\ell$}
{$4\ell$}
{\True $3\ell$}
{$1\ell$}
\loigiai{
Lượng nước múc ra: $250 \times 8 = 2000\text{ ml} = 2\ell$. Trong bể còn lại: $5 - 2 = 3\ell$.
\textbf{Đáp án của câu hỏi là C}
}
\end{ex}

\textbf{Phần 2. Trắc nghiệm đúng sai}
\setcounterpageref{ex}{0}

\begin{ex} %[3D3V32]
Xét tính đúng sai của các câu nói về mi-li-lít:
\choiceTF
{\True Mi-li-lít được viết tắt là $\text{ml}$.}
{$1\ell = 100\text{ ml}$.}
{\True $500\text{ ml}$ bằng nửa lít nước.}
{$750\text{ ml} + 350\text{ ml} = 1\ell$.}
\loigiai{
\begin{itemize}
    \item Ý a) Đúng: Kí hiệu quy chuẩn.
    \item Ý b) Sai: Sai vì $1\ell = 1000\text{ ml}$.
    \item Ý c) Đúng: $1000 : 2 = 500\text{ ml}$.
    \item Ý d) Sai: Sai vì tổng bằng $1100\text{ ml} = 1\ell\text{ }100\text{ ml}$.
\end{itemize}
\begin{itemize}
    \item a) \textbf{Đúng}.
    \item b) \textbf{Sai}.
    \item c) \textbf{Đúng}.
    \item d) \textbf{Sai}.
\end{itemize}
}
\end{ex}

\begin{ex} %[3D3V32]
Một ấm đun nước chứa đầy $1\ell\text{ }500\text{ ml}$ nước:
\choiceTF
{\True Ấm chứa tất cả $1500\text{ ml}$ nước.}
{\True Nếu rót ra 3 cốc, mỗi cốc $300\text{ ml}$ thì đã rót ra $900\text{ ml}$.}
{Sau khi rót ra $900\text{ ml}$, trong ấm còn lại $700\text{ ml}$.}
{\True Lượng nước còn lại ($600\text{ ml}$) đủ rót thêm 2 cốc $300\text{ ml}$ nữa.}
\loigiai{
\begin{itemize}
    \item Ý a) Đúng: Đổi $1\ell\text{ }500\text{ ml} = 1500\text{ ml}$.
    \item Ý b) Đúng: $300 \times 3 = 900\text{ ml}$.
    \item Ý c) Sai: Sai vì còn lại $1500 - 900 = 600\text{ ml}$.
    \item Ý d) Đúng: $600 : 300 = 2$ cốc.
\end{itemize}
\begin{itemize}
    \item a) \textbf{Đúng}.
    \item b) \textbf{Đúng}.
    \item c) \textbf{Sai}.
    \item d) \textbf{Đúng}.
\end{itemize}
}
\end{ex}

\textbf{Phần 3. Trả lời ngắn}
\setcounterpageref{ex}{0}

\begin{ex} %[3D3V32]
Đổi $1\ell - 400\text{ ml}$ ra đơn vị mi-li-lít.
\shortans{600}
\loigiai{
\textbf{Đáp số:} $600$\\
$1000 - 400 = 600\text{ (ml)}$.
}
\end{ex}

\begin{ex} %[3D3V32]
Một lọ mật ong có dung tích $350\text{ ml}$. Hỏi 2 lọ mật ong như thế có tổng dung tích là bao nhiêu mi-li-lít?
\shortans{700}
\loigiai{
\textbf{Đáp số:} $700$\\
$350 \times 2 = 700\text{ (ml)}$.
}
\end{ex}

\begin{ex} %[3D3V32]
Người ta chia đều $800\text{ ml}$ nước trái cây vào 4 chiếc ly. Hỏi mỗi ly chứa bao nhiêu mi-li-lít?
\shortans{200}
\loigiai{
\textbf{Đáp số:} $200$\\
$800 : 4 = 200\text{ (ml)}$.
}
\end{ex}

\begin{ex} %[3D3V32]
Cần bao nhiêu hộp sữa tươi loại $250\text{ ml}$ để có được đúng $1\ell$ sữa?
\shortans{4}
\loigiai{
\textbf{Đáp số:} $4$\\
$1000 : 250 = 4$ hộp.
}
\end{ex}

\begin{ex} %[3D3V32]
Một bình nước có $900\text{ ml}$. Bạn An uống $\dfrac{1}{3}$ lượng nước đó. Hỏi trong bình còn lại bao nhiêu mi-li-lít nước?
\shortans{600}
\loigiai{
\textbf{Đáp số:} $600$\\
An uống: $900 : 3 = 300\text{ ml}$. Còn lại: $900 - 300 = 600\text{ (ml)}$.
}
\end{ex}

\begin{ex} %[3D3V32]
Tính giá trị của biểu thức: $120\text{ ml} \times 5 - 150\text{ ml}$ theo đơn vị mi-li-lít.
\shortans{450}
\loigiai{
\textbf{Đáp số:} $450$\\
$600 - 150 = 450\text{ (ml)}$.
}
\end{ex}
